% Nouveau EGA V (cote 25): archive cote 25, site page 95 (PDF page 96).
% typescript. Draft compared: none.
% AI-transcribed by Claude Opus 5.5 (claude-opus-5-5, adaptive thinking, effort high), Message Batches API, batch ega25b (msgbatch_01A7n5YuXubUPNK1zsepi8RG), 2026-11-13; not yet checked by eye.
\scanpage{25/95}
\mspage{89}

Donc dim E \hand{$<$} N = dim $P$ , donc $E_{\hand{$y$}} = \emptyset$ , et à fortiori $X_{\hand{$y$}}$ est lisse\hand{.}\inserted{\unsure{sur $k(y)$}}
Ce point acquis, on conclut aussitôt à la validité des variantes
évidentes de 2.12. à 2.18. dans le contexte actuel.

Le passage au nº4 d'une section générique à une section générale
et les développements du nº5, sont évidemment valables dans le contexte
actuel (mais sont à tel point des tautologies ou redites\xout{\illegible{1}} des
par.8,9,12 \xout{)} qu'on hésite à les énoncer en forme). De même les
développements de 7.1. , valables en tous cas si k est \xout{\unsure{infini}} alg clos
(et même si k est simplement infini, si on suppose G rationelle\xout{s} sur k),
et les cas particuliers 7.2. , 7.3. ; quant au\xout{x} resultat\xout{s} 7.4. \xout{\unsure{et 7.5.}}
\xout{\illegible{1}}c'est évidemment \inserted[typed]{une application} de nature spéciale \inserted[typed]{de} \xout{à} la situation des sections
hyperplanes . Comme j'ai dit, les nºs 3 et 6 sont suspendus à l'exten-%
sion du th. de Zariski.

Il resterait à étendre également les résultats du nº8 (respris au
nº 12), qui prendront peut-être comme-ça une allure plus sympathique.
Je te conseille même de commencer par rédiger ces résultats dans ce cadre,
en essayant de pousser aussi loin que possible dans cette voie. J'ai
l'impression qu'on doit pouvoir récupérer au moins tout ce qui n'est pas
conséquence directe de 8.7. c) (et encore pourrait-on essayer
d'abstraire des conditions axiomatiques faisant marcher une variante de
8.7. c). Je me borne à ces recommandations, mais suis prêt à y revenir
avec plus de détails si tu as des difficultés particulières.

\mspage{95}
